% =====================================================================
\section{5.2 소프트 마진: 완벽히 나뉘지 않는 데이터}
\begin{frame}{도입: 하드 마진의 치명적인 전제}
  \begin{itemize}
    \item 5.1절 하드 마진: 제약 $y^{(i)}(w^\top x^{(i)}+b)\ge 1$ 아래
      $\|w\|$ 최소화. 쌍대 문제, 서포트 벡터, $2/\|w\|$ 공식 등
      예쁜 성질을 다 갖췄다.
    \item 그런데 \textbf{치명적인 전제}: \kb{데이터가 완벽히 선형분리 가능}.
  \end{itemize}
  \vskip0.5em
  실전 데이터는 노이즈 때문에 이 가정을 거의 만족하지 못한다:
  \begin{itemize}
    \item \kb{라벨링 오류} — 의사가 진단을 틀렸거나,
      스팸 필터 학습용 메일에 사람이 틀린 태그를 붙인 경우
    \item \kb{측정 오류} — 센서 오차
    \item \kb{모호한 경계 사례} — 도메인 밖의 데이터
  \end{itemize}
  \vskip0.4em
  \begin{block}{문제의 심각성}
    어떤 $(w,b)$로도 제약을 모두 만족 못 하면(불일치, infeasible)
    이차계획을 풀 수 없다 — 하드 마진 SVM은 \textbf{라벨 오류
    단 하나만 있어도 해 자체가 존재하지 않는} 모델. 그대로 실전에
    쓸 수 없다.
  \end{block}
\end{frame}
